\newpage
\phantomsection
\label{prf:density-of-rationals-in-reals}
\LRAProofFor{thm:density-of-rationals-in-reals}

\begin{remark*}[Return]
\hyperref[thm:density-of-rationals-in-reals]{Return to Theorem}
\end{remark*}

\begin{theorem*}[Density of the Rationals]
For every $a,b\in\mathbb{R}$ with $a<b$, there exists $q\in\mathbb{Q}$ such that $a<q<b$.
\end{theorem*}

\begin{proof}[Professional Standard Proof]
\LRAProofBodyStart
Let $a,b\in\mathbb{R}$ with $a<b$. Then $b-a>0$.
By the Archimedean reciprocal form, choose $N\in\mathbb{N}$ such that
\[
0<\frac{1}{N}<b-a.
\]
Equivalently,
\[
1<N(b-a)=Nb-Na.
\]
Hence the interval $(Na,Nb)$ has length greater than $1$. By the unit-length
interval lemma, there exists $k\in\mathbb{Z}$ such that
\[
Na<k<Nb.
\]
Since $N>0$, division by $N$ preserves the inequalities, so
\[
a<\frac{k}{N}<b.
\]
Because $k\in\mathbb{Z}$ and $N\in\mathbb{N}$ with $N\ne 0$, the number
$k/N$ is rational. Thus $q=k/N\in\mathbb{Q}$ satisfies $a<q<b$.
\end{proof}

\begin{proof}[Detailed Learning Proof]
\LRAProofBodyStart
Let $a,b\in\mathbb{R}$ and suppose $a<b$. The goal is to construct a rational
number strictly between them.

Since $a<b$, the distance between the endpoints is positive:
\[
b-a>0.
\]
Apply the Archimedean reciprocal form to the positive number $b-a$. There is a
natural number $N\in\mathbb{N}$ such that
\[
0<\frac{1}{N}<b-a.
\]
Multiplying the strict inequality $1/N<b-a$ by the positive number $N$ gives
\[
1<N(b-a).
\]
That is,
\[
Nb-Na>1.
\]
So the scaled interval
\[
(Na,Nb)
\]
has length greater than one.

An open interval of length greater than one contains an integer. Therefore
there exists $k\in\mathbb{Z}$ such that
\[
Na<k<Nb.
\]
Now scale back down by dividing through by the positive natural number $N$.
Because $N>0$, the direction of the inequalities is unchanged:
\[
a<\frac{k}{N}<b.
\]
Finally, $k/N$ is rational because it is a quotient of an integer by a nonzero
integer. Setting
\[
q=\frac{k}{N}
\]
gives $q\in\mathbb{Q}$ and $a<q<b$. Hence every nonempty open interval in
$\mathbb{R}$ contains a rational number.
\end{proof}

\begin{remark*}[Proof structure]
First choose a scale $N$ so fine that $1/N<b-a$. Then the scaled interval
$(Na,Nb)$ has length greater than one, so it contains an integer $k$. Scaling
that integer back down produces the rational number $k/N$ inside $(a,b)$.
\end{remark*}

\begin{dependencies}
\hyperref[cor:archimedean-reciprocal-form]{Archimedean Reciprocal Form},
\hyperref[lem:unit-length-interval-contains-integer]{Unit-Length Interval Contains an Integer}.
\end{dependencies}

\clearpage
